% ============================================================================
\chapter{ramfs：内存文件系统}
\label{ch:ramfs}
% Stage E-2 ~ E-3
% ============================================================================

\section{本章目标}

\begin{itemize}
  \item 实现 ramfs 后端：内存中的简易文件系统
  \item 实现 initrd 加载：GRUB Multiboot2 module → ramfs 根目录
  \item ELF 加载器通过 VFS 读取文件，而非硬编码内存地址
  \item 验证 \texttt{exec("/hello")} 从 initrd 加载 ELF 并运行
\end{itemize}

% ============================================================================
\section{概念：initrd 与内存文件系统}
\label{sec:ramfs-concept}
% ============================================================================

% TODO:
% - 鸡生蛋问题：磁盘驱动还没实现，用户程序放在哪里？
%     答案：打包进内核镜像（initrd = initial ramdisk）
% - Linux 方案：initramfs（CPIO 格式的压缩包，解压到 ramfs）
% - 我们的方案：GRUB 以 Multiboot2 module 加载 initrd → 解析到 ramfs
% - ramfs = 纯内存文件系统：文件内容和目录结构都在 RAM 中
%     优点：简单快速，无需磁盘驱动
%     缺点：断电即丢失，大小受限于内存

\subsection{打包格式选择}

% TODO:
% - 方案对比：
%     CPIO（newc 格式）：Linux 标准，工具支持好，解析简单
%     tar：通用但解析稍复杂（header 512 bytes 对齐）
%     自定义格式：最简单但不通用
% - 推荐 CPIO newc：定长 ASCII header，字段含义清晰
% - CPIO header 结构图

% ============================================================================
\section{ramfs 后端实现}
\label{sec:ramfs-impl}
% ============================================================================

% TODO:
% - 内部数据结构：
%     ramfs_node { name, type(FILE/DIR), data_ptr, data_size, children_list }
% - 实现 fs_ops_t 全部方法：
%     ramfs_mount：初始化根节点
%     ramfs_open：路径查找 → 返回 file
%     ramfs_read：从 data_ptr + offset 复制到 buf
%     ramfs_write：扩展 data（kmalloc 重新分配）
%     ramfs_stat：返回 size/mode
%     ramfs_readdir：遍历 children_list
% - ramfs_get_ops() 返回 fs_ops_t 指针 → 注册到 VFS

% ============================================================================
\section{initrd 加载}
\label{sec:ramfs-initrd}
% ============================================================================

% TODO:
% - Multiboot2 module tag 解析：
%     tag type 3 (MODULE): mod_start, mod_end, cmdline
%     物理地址 → PHYS_TO_VIRT 转换
% - CPIO 解析器：
%     遍历 CPIO entries → 对每个文件调用 ramfs_create(path, data, size)
%     遇到 "TRAILER!!!" 结束
% - grub.cfg 配置：
%     module2 /boot/initrd.cpio
% - 构建脚本：
%     find user_programs/ | cpio -o -H newc > initrd.cpio

% ============================================================================
\section{验证}
\label{sec:ramfs-verify}
% ============================================================================

% TODO:
% - 测试 1：ramfs 挂载 → readdir("/") → 看到 initrd 中的文件列表
% - 测试 2：open("/hello") → read 完整内容 → 与原文件对比
% - 测试 3：exec("/hello") → 通过 VFS 读取 ELF → 加载执行 → 串口输出
% - shell 命令：ls / cat / fs test

% ============================================================================
\section{本章小结}
% ============================================================================

ramfs 是 \texttt{fs\_ops\_t} 的第一个具体后端，证明了 VFS 可插拔架构的可行性。
配合 initrd，用户程序无需磁盘驱动即可被加载执行。
下一章实现块设备抽象和磁盘文件系统，让内核能访问持久化存储。
